\documentclass[11pt,a4paper]{article}

\usepackage{amsmath, amssymb, amsthm}
\usepackage{physics}
\usepackage{bm}
\usepackage{siunitx}
\usepackage{booktabs}
\usepackage{graphicx}
\usepackage[margin=25mm]{geometry}
\usepackage{hyperref}

\title{Technical Note: block-DDA\_Py\\[4pt]
\large Discrete Dipole Approximation with FFT-accelerated\\
Block-Krylov Solver for Light Scattering by Arbitrary Particles}
\author{N.\ Moteki}
\date{Last updated: 2026-03-26 (corresponds to block-DDA\_Py v0.6.0)}

\begin{document}
\maketitle
\tableofcontents
\newpage

%======================================================================
\section{Introduction}
\label{sec:introduction}
%======================================================================

The discrete dipole approximation (DDA) is a flexible numerical method for
computing electromagnetic scattering and absorption by particles of arbitrary
shape and composition~\cite{Purcell1973,Draine1988}.
The target particle is represented by an array of $N$ point dipoles on a cubic
lattice with spacing~$d$.  Each dipole acquires a polarisation in response to
the local electric field, which consists of the incident wave plus the fields
radiated by all other dipoles.

This document describes the theoretical background and the algorithms
that are implemented in version~0.4.0 of
\texttt{block-DDA\_Py}, covering:
\begin{enumerate}
  \item the DDA formulation and polarisability prescription
        (Sec.~\ref{sec:dda-formulation}),
  \item the Goodman~(1991) FFT-accelerated matrix--vector product
        (Sec.~\ref{sec:goodman-fft}),
  \item the block-Krylov iterative solvers
        (Sec.~\ref{sec:block-krylov}),
  \item the scattering observables computed from the solution
        (Sec.~\ref{sec:observables}),
  \item the Gaussian Random Ellipsoid (GRE) shape model
        (Sec.~\ref{sec:gre}),
  \item the Mie theory reference solution for homogeneous spheres
        (Sec.~\ref{sec:mie}).
\end{enumerate}

Throughout this note, the electromagnetic equations are written in
Gaussian (CGS) units; lengths are expressed in micrometers (\si{\micro\metre}).
Vectors are set in bold italic ($\bm{E}$, $\bm{k}$, $\bm{r}$),
matrices and second-rank tensors in upright bold ($\mathbf{A}$,
$\boldsymbol{\alpha}$), and unit vectors carry a hat
($\hat{\bm{e}}$, $\hat{\bm{r}}$).

%======================================================================
\section{DDA Formulation}
\label{sec:dda-formulation}
%======================================================================

%----------------------------------------------------------------------
\subsection{Fundamental equation}
\label{sec:fundamental-eq}
%----------------------------------------------------------------------

Consider a monochromatic plane wave of vacuum wavelength $\lambda_0$ incident
on a particle embedded in a non-absorbing medium of refractive index~$m_m$.
The wavenumber in the medium is
\begin{equation}
  k = \frac{2\pi\, m_m}{\lambda_0}.
  \label{eq:wavenumber}
\end{equation}

The particle is discretised into $N$ cubic elements (dipoles) of side
length~$d$ and volume $v = d^3$, placed on a regular lattice with grid
dimensions $N_x \times N_y \times N_z$.
The lattice position of the $i$-th dipole is $\bm{r}_i$.

The polarisation $\bm{P}_i$ of dipole~$i$ satisfies
\begin{equation}
  \bm{P}_i = \boldsymbol{\alpha}_i \cdot \bm{E}(\bm{r}_i),
  \label{eq:polarisation}
\end{equation}
where $\boldsymbol{\alpha}_i$ is the $3\times 3$ polarisability tensor and
$\bm{E}(\bm{r}_i)$ is the total electric field at position~$\bm{r}_i$.
The total field is the superposition of the incident field and the
contributions from all other dipoles:
\begin{equation}
  \bm{E}(\bm{r}_i)
    = \bm{E}_{\text{inc},i}
      - \sum_{j \neq i} \mathbf{G}(\bm{r}_i - \bm{r}_j) \cdot \bm{P}_j,
  \label{eq:total-field}
\end{equation}
where $\mathbf{G}$ is the free-space dyadic Green's function
(Sec.~\ref{sec:green-tensor}).
Substituting \eqref{eq:total-field} into \eqref{eq:polarisation} and
rearranging yields the $3N \times 3N$ linear system
\begin{equation}
  \mathbf{A} \, \bm{P} = \bm{E}_{\text{inc}},
  \label{eq:dda-system}
\end{equation}
where $\bm{P}$ and $\bm{E}_{\text{inc}}$ are column vectors of length~$3N$
formed by concatenating the three Cartesian components of each dipole
(ordering: $P_{x,0}, P_{y,0}, P_{z,0}, P_{x,1}, \ldots$),
and $\mathbf{A}$ is a complex-valued matrix with
\begin{equation}
  \mathbf{A}_{ij} =
  \begin{cases}
    \boldsymbol{\alpha}_i^{-1}         & \text{if } i = j, \\[4pt]
    \mathbf{G}(\bm{r}_i - \bm{r}_j)   & \text{if } i \neq j.
  \end{cases}
  \label{eq:A-matrix}
\end{equation}

When the dipoles are placed on a regular cubic lattice, the off-diagonal
blocks $\mathbf{A}_{ij}$ depend only on the difference $\bm{r}_i - \bm{r}_j$,
giving $\mathbf{A}$ a block-Toeplitz structure that can be exploited by FFT
(Sec.~\ref{sec:goodman-fft}).

%----------------------------------------------------------------------
\subsection{Incident field}
\label{sec:incident-field}
%----------------------------------------------------------------------

The incident plane wave in the medium is
\begin{equation}
  \bm{E}_{\text{inc},i}
    = \bm{E}_0 \exp(i \bm{k} \cdot \bm{r}_i - i\omega t),
  \label{eq:incident-wave}
\end{equation}
where $\bm{k} = k\,\hat{\bm{u}}_{\text{inc}}$ with
$\hat{\bm{u}}_{\text{inc}}$ the propagation direction
(fixed to $+z$ in the laboratory frame).
The time-harmonic factor $e^{-i\omega t}$ is suppressed hereafter.

To simulate orientation-averaged scattering, the particle is rotated through
$L = N_\alpha \times N_\beta \times N_\gamma$ orientations on a
deterministic grid of intrinsic ZYZ Euler angles:
\begin{itemize}
  \item $\alpha$: $N_\alpha$ equally spaced values in $[0, 2\pi)$,
  \item $\beta$: $N_\beta$ values chosen so that $\cos\beta$ is equally
        spaced in $(-1, 1)$ (equal-area midpoints on the sphere),
  \item $\gamma$: $N_\gamma$ equally spaced values in $[0, 2\pi)$.
\end{itemize}
This deterministic scheme avoids sampling noise and ensures reproducibility.
For each orientation, the laboratory-frame vectors are transformed into
the particle frame via the inverse rotation matrix.
Batch rotation is performed using \texttt{scipy.spatial.transform.Rotation}.

Following the CAS-v2 (Complex Amplitude Sensing, version~2) beam
configuration, the incident polarisation is assumed to be left-handed
circular.
This setting is hard-coded in the current version (v0.6.0).
The polarisation vector is:
\begin{equation}
  \hat{\bm{e}}_0
    = \frac{1}{\sqrt{2}}\,\hat{\bm{e}}_\theta
      + \frac{i}{\sqrt{2}}\,\hat{\bm{e}}_\phi,
  \label{eq:circular-pol}
\end{equation}
where $\hat{\bm{e}}_\theta$ and $\hat{\bm{e}}_\phi$ are the polar and
azimuthal unit vectors associated with the incident direction.

%----------------------------------------------------------------------
\subsection{Free-space dyadic Green's function}
\label{sec:green-tensor}
%----------------------------------------------------------------------

The $3\times 3$ interaction tensor for two dipoles separated by
$\bm{r} = \bm{r}_i - \bm{r}_j$ ($r = |\bm{r}|$, $\hat{\bm{r}} = \bm{r}/r$)
is~\cite{Goodman1991}
\begin{equation}
  \mathbf{G}(\bm{r})
    = \frac{e^{ikr}}{r}
      \left[
        k^2 \bigl(\hat{\bm{r}} \otimes \hat{\bm{r}} - \mathbf{I}_3\bigr)
        + \frac{1 - ikr}{r^2}
          \bigl(\mathbf{I}_3 - 3\,\hat{\bm{r}} \otimes \hat{\bm{r}}\bigr)
      \right],
  \label{eq:green-tensor}
\end{equation}
where $\otimes$ denotes the outer (tensor) product and $\mathbf{I}_3$ the
$3\times 3$ identity matrix.
The self-interaction ($\bm{r} = \bm{0}$) is excluded from
$\mathbf{G}$ and handled separately via the diagonal
$\boldsymbol{\alpha}_i^{-1}$ in \eqref{eq:A-matrix}.

%----------------------------------------------------------------------
\subsection{Polarisability: Clausius--Mossotti with radiative correction}
\label{sec:polarisability}
%----------------------------------------------------------------------

The relative permittivity of the $i$-th dipole along Cartesian axis~$c$
($c \in \{x, y, z\}$) is
\begin{equation}
  \epsilon_{r,c} = \frac{m_{p,c}^2}{m_m^2},
  \label{eq:eper-r}
\end{equation}
where $m_{p,c}$ is the (generally complex) particle refractive index along
axis~$c$.

The static Clausius--Mossotti (CM) polarisability is
\begin{equation}
  \alpha_{0,c}
    = \frac{3}{4\pi}\,
      \frac{\epsilon_{r,c} - 1}{\epsilon_{r,c} + 2}\; v,
  \label{eq:cm-polarisability}
\end{equation}
where $v = d^3$ is the dipole volume.

To improve accuracy, the CM polarisability is corrected for radiative
reaction following Chaumet \& Rahmani~\cite{Chaumet2009}.
Define the equivalent-volume sphere radius
\begin{equation}
  a = \left(\frac{3v}{4\pi}\right)^{\!1/3}
  \label{eq:equiv-radius-a}
\end{equation}
and the correction term
\begin{equation}
  M = \frac{8\pi}{3}
      \bigl[(1 - ika)\,e^{ika} - 1\bigr].
  \label{eq:M-term}
\end{equation}
The dynamical (radiative-reaction corrected) polarisability is then
\begin{equation}
  \alpha_c
    = \frac{\alpha_{0,c}}{1 - M\,\alpha_{0,c}\,/\,v}.
  \label{eq:rr-polarisability}
\end{equation}

In general the particle may be optically anisotropic, so
$\boldsymbol{\alpha}_i = \operatorname{diag}(\alpha_x, \alpha_y, \alpha_z)$.
For isotropic particles, $m_{p,x} = m_{p,y} = m_{p,z}$ and the
polarisability reduces to a scalar times $\mathbf{I}_3$.

%======================================================================
\section{FFT-Accelerated Matrix--Vector Product (Goodman 1991)}
\label{sec:goodman-fft}
%======================================================================

Direct evaluation of the matrix--vector product (MVP)
$\mathbf{A}\,\bm{P}$ requires $O(N^2)$ operations and $O(N^2)$ storage.
Because the off-diagonal part of $\mathbf{A}$ has block-Toeplitz structure
on a regular cubic lattice, the MVP can be evaluated via 3D FFT in
$O(N \log N)$ time with $O(N)$ storage~\cite{Goodman1991}.

%----------------------------------------------------------------------
\subsection{Interaction tensor on the doubled grid}
\label{sec:doubled-grid}
%----------------------------------------------------------------------

The physical cuboid lattice has dimensions $N_x \times N_y \times N_z$.
The index difference along each axis~$\alpha$ ranges from
$-(N_\alpha - 1)$ to $+(N_\alpha - 1)$.
To evaluate the linear convolution as a cyclic convolution via FFT,
the lattice is extended to $2N_x \times 2N_y \times 2N_z$ with the
following index mapping.

For each extended-grid index $I_\alpha \in \{0, 1, \ldots, 2N_\alpha - 1\}$,
the physical integer displacement $\Delta\alpha$ is
\begin{equation}
  \Delta\alpha(I_\alpha) =
  \begin{cases}
    I_\alpha            & 0 \le I_\alpha \le N_\alpha - 1
                          \quad\text{(positive displacements)}, \\
    0                   & I_\alpha = N_\alpha
                          \quad\text{(Nyquist; zeroed out)}, \\
    I_\alpha - 2N_\alpha & N_\alpha + 1 \le I_\alpha \le 2N_\alpha - 1
                          \quad\text{(negative displacements)}.
  \end{cases}
  \label{eq:index-mapping}
\end{equation}

The physical displacement vector at extended-grid index
$(I_x, I_y, I_z)$ is
$\bm{r} = d\,(\Delta x,\, \Delta y,\, \Delta z)$,
and the interaction tensor $\mathbf{G}(\bm{r})$ is computed via
\eqref{eq:green-tensor}.
Three special cases are set to zero:
\begin{itemize}
  \item self-interaction at the origin:
        $\mathbf{G}(\bm{0}) = \mathbf{0}$,
  \item Nyquist positions ($I_\alpha = N_\alpha$ for any~$\alpha$):
        physically unreachable displacements.
\end{itemize}

The resulting array has shape
$(2N_x, 2N_y, 2N_z, 3, 3)$ and dtype \texttt{complex128}.

%----------------------------------------------------------------------
\subsection{Pre-computation of the FFT}
\label{sec:fft-precompute}
%----------------------------------------------------------------------

Before the iterative solver begins, the 3D discrete Fourier transform of
$\mathbf{G}$ along the spatial axes is computed once:
\begin{equation}
  \widetilde{\mathbf{G}}
    = \mathrm{FFT}_{3\text{D}}\!\bigl[\mathbf{G}\bigr].
  \label{eq:fft-precompute}
\end{equation}
$\widetilde{\mathbf{G}}$ has the same shape
$(2N_x, 2N_y, 2N_z, 3, 3)$ and is stored in memory throughout the
simulation.
It is independent of the incident direction, polarisation, and Krylov
iteration, so this computation is performed only once per
$(\lambda_0, m_m)$ pair.

%----------------------------------------------------------------------
\subsection{Block MVP via FFT convolution}
\label{sec:fft-mvp}
%----------------------------------------------------------------------

Given a block of $L$ polarisation vectors stacked as columns in
$\mathbf{P}_{\text{block}}$ (shape $3N_xN_yN_z \times L$), the
off-diagonal MVP $\mathbf{A}_{\text{off}} \, \mathbf{P}_{\text{block}}$
is computed as follows.

\begin{enumerate}
  \item \textbf{Reshape} $\mathbf{P}_{\text{block}}$ from
        $(3 N_x N_y N_z,\, L)$ to the 5D physical representation
        $(N_x, N_y, N_z, 3, L)$.
  \item \textbf{Zero-pad} to the doubled grid:
        \begin{equation}
          \hat{\mathbf{P}}(I_x, I_y, I_z, c, \ell) =
          \begin{cases}
            \mathbf{P}(I_x, I_y, I_z, c, \ell)
              & \text{if } I_\alpha < N_\alpha \;\forall \alpha, \\
            0 & \text{otherwise}.
          \end{cases}
          \label{eq:zero-pad}
        \end{equation}
  \item \textbf{Forward 3D FFT} along spatial axes:
        $\widetilde{\mathbf{P}} = \mathrm{FFT}_{3\text{D}}[\hat{\mathbf{P}}]$.
  \item \textbf{Pointwise tensor product} in Fourier space:
        \begin{equation}
          \widetilde{Y}_{i,\ell}
            = \sum_{j=1}^{3}
              \widetilde{G}_{ij}\;\widetilde{P}_{j,\ell}
          \qquad
          \text{(at each spatial frequency)},
          \label{eq:fourier-product}
        \end{equation}
        i.e.\ a $3\times 3$ matrix times a $3\times L$ matrix at every
        grid point.
  \item \textbf{Inverse 3D FFT}:
        $\hat{\mathbf{Y}} = \mathrm{FFT}_{3\text{D}}^{-1}[\widetilde{\mathbf{Y}}]$.
  \item \textbf{Extract} the physical region
        $\hat{\mathbf{Y}}[{:}N_x,\, {:}N_y,\, {:}N_z,\, :,\, :]$
        and flatten back to $(3 N_x N_y N_z, L)$.
\end{enumerate}

The full preconditioned MVP including the diagonal is
\begin{equation}
  \mathbf{A}\,\bm{S}
    = \operatorname{diag}(\mathbf{A})\cdot\bm{S}
      + \mathbf{A}_{\text{off}}\,\bm{S},
  \label{eq:full-mvp}
\end{equation}
where $\operatorname{diag}(\mathbf{A}) = \boldsymbol{\alpha}^{-1}$ is applied
element-wise and $\mathbf{A}_{\text{off}}$ is evaluated via the FFT
procedure above.

%----------------------------------------------------------------------
\subsection{Computational complexity and memory}
\label{sec:complexity}
%----------------------------------------------------------------------

\begin{itemize}
  \item \textbf{Time}: Each MVP costs $O(N_{\text{cuboid}} \log N_{\text{cuboid}})$
        where $N_{\text{cuboid}} = N_x N_y N_z$.
        For the block variant with $L$ right-hand sides, the $L$ columns
        are processed in a single batched FFT call.
  \item \textbf{Memory}: The dominant allocations are:
  \begin{itemize}
    \item $\widetilde{\mathbf{G}}$:
          $2N_x \times 2N_y \times 2N_z \times 9$ complex128 entries
          $= 1152\,N_{\text{cuboid}}$ bytes (static, allocated once);
    \item padded polarisation arrays (forward and inverse FFT):
          $2N_x \times 2N_y \times 2N_z \times 3 \times L$ complex128
          entries each $= 384\,L\,N_{\text{cuboid}}$ bytes (two such arrays
          at peak).
  \end{itemize}
  Peak memory $\approx (1152 + 768\,L)\,N_{\text{cuboid}}$ bytes.
\end{itemize}

%----------------------------------------------------------------------
\subsection{Multi-threaded FFT}
\label{sec:fft-workers}
%----------------------------------------------------------------------

The 3D FFTs are computed by \texttt{scipy.fft.fftn} with multi-threading.
The default number of worker threads is
$\max(1,\; n_{\text{cpu}} - 2)$, leaving two cores free for other
processes.
This can be overridden at runtime via the environment variable
\texttt{DDA\_FFT\_WORKERS}.

%======================================================================
\section{Block-Krylov Iterative Solvers}
\label{sec:block-krylov}
%======================================================================

The DDA system \eqref{eq:dda-system} is solved for $L$ right-hand sides
simultaneously.
Stacking the $L$ incident-field vectors as columns of a block matrix
$\mathbf{B}$ (shape $3N \times L$), the problem becomes
\begin{equation}
  \mathbf{A}\,\mathbf{X} = \mathbf{B},
  \label{eq:block-system}
\end{equation}
where $\mathbf{X}$ (shape $3N \times L$) contains the $L$ solution
vectors.

%----------------------------------------------------------------------
\subsection{Jacobi preconditioning}
\label{sec:jacobi}
%----------------------------------------------------------------------

A left Jacobi (diagonal) preconditioner is applied:
\begin{equation}
  \mathbf{M}^{-1} = \operatorname{diag}(\mathbf{A})^{-1}
                   = \operatorname{diag}(\alpha_{c,i}),
  \label{eq:jacobi-precond}
\end{equation}
so that the preconditioned system becomes
$\mathbf{M}^{-1}\mathbf{A}\,\mathbf{X} = \mathbf{M}^{-1}\mathbf{B}$.

%----------------------------------------------------------------------
\subsection{Block-BiCGSTAB}
\label{sec:bicgstab}
%----------------------------------------------------------------------

The primary solver is Block-BiCGSTAB~\cite{Tadano2009}, which extends
the stabilised bi-conjugate gradient method to block systems.
It is suitable for general (non-symmetric) matrices and is the default
in \texttt{block-DDA\_Py}.

The algorithm is summarised in Table~\ref{tab:bicgstab}.
Each iteration requires two block MVPs (lines~1 and~4).
The default convergence parameters are $\tau = 0.01$ and
$n_{\max} = 25$.

\begin{table}[htbp]
  \caption{Block-BiCGSTAB with Jacobi preconditioning
           (after Tadano et al.~\cite{Tadano2009}).}
  \label{tab:bicgstab}
  \centering
  \small
  \renewcommand{\arraystretch}{1.4}
  \begin{tabular}{@{}r@{\quad}l@{}}
    \toprule
    \multicolumn{2}{@{}l}{\textbf{Input:}
      $\mathbf{A}$, $\mathbf{B}$ (shape $3N\!\times\!L$),
      $\mathbf{M}=\operatorname{diag}(\mathbf{A})$,
      tolerance $\tau$, max iterations $n_{\max}$.} \\
    \midrule
    \multicolumn{2}{@{}l}{\textbf{Initialise:}\;
      $\mathbf{X}_0 = \mathbf{0}$,\;
      $\mathbf{R}_0 = \mathbf{M}^{-1}\mathbf{B}$,\;
      $\mathbf{P}_0 = \mathbf{R}_0$,\;
      $\widetilde{\mathbf{R}}_0 = \mathbf{R}_0$.} \\
    \midrule
    \multicolumn{2}{@{}l}{\textbf{For} $n = 0, 1, \ldots, n_{\max}-1$:} \\[2pt]
    1. & $\mathbf{V}_n
          = \mathbf{M}^{-1}\mathbf{A}\,\mathbf{P}_n$ \\
    2. & $\boldsymbol{\alpha}_n
          = (\widetilde{\mathbf{R}}_0^\mathsf{H}\,\mathbf{V}_n)^{-1}\;
            \widetilde{\mathbf{R}}_0^\mathsf{H}\,\mathbf{R}_n$ \\
    3. & $\mathbf{T}_n
          = \mathbf{R}_n - \mathbf{V}_n\,\boldsymbol{\alpha}_n$ \\
    4. & $\mathbf{Z}_n
          = \mathbf{M}^{-1}\mathbf{A}\,\mathbf{T}_n$ \\
    5. & $\xi_n
          = \operatorname{tr}(\mathbf{Z}_n^\mathsf{H}\,\mathbf{T}_n)\;/\;
            \operatorname{tr}(\mathbf{Z}_n^\mathsf{H}\,\mathbf{Z}_n)$ \\
    6. & $\mathbf{X}_{n+1}
          = \mathbf{X}_n + \mathbf{P}_n\,\boldsymbol{\alpha}_n
            + \xi_n\,\mathbf{T}_n$ \\
    7. & $\mathbf{R}_{n+1}
          = \mathbf{T}_n - \xi_n\,\mathbf{Z}_n$ \\
    8. & If $\|\mathbf{R}_{n+1}\| / \|\mathbf{M}^{-1}\mathbf{B}\|
          < \tau$: stop, return $\mathbf{X}_{n+1}$. \\
    9. & $\boldsymbol{\beta}_n
          = (\widetilde{\mathbf{R}}_0^\mathsf{H}\,\mathbf{V}_n)^{-1}\;
            (-\widetilde{\mathbf{R}}_0^\mathsf{H}\,\mathbf{Z}_n)$ \\
    10. & $\mathbf{P}_{n+1}
           = \mathbf{R}_{n+1}
             + (\mathbf{P}_n - \xi_n\,\mathbf{V}_n)\,
               \boldsymbol{\beta}_n$ \\
    \bottomrule
  \end{tabular}
\end{table}

%----------------------------------------------------------------------
\subsection{Block-COCG-rq (alternative solver)}
\label{sec:cocg}
%----------------------------------------------------------------------

Since the DDA interaction matrix $\mathbf{A}$ is complex symmetric
($\mathbf{A} = \mathbf{A}^\top$), a solver specialised for this
structure, Block-COCG-rq~\cite{Gu2016}, is also provided.
Block-COCG-rq requires only one block MVP per iteration and is
therefore faster per iteration than Block-BiCGSTAB.
However, for particles with high refractive index it exhibits
poorer convergence stability compared to Block-BiCGSTAB.
To switch the solver, replace the import and function call in
\texttt{bl\_dda/scatterer.py} (line~3 and line~169):
change \texttt{bl\_bicgstab\_jacobi\_mvp\_fft} to
\texttt{bl\_cocg\_rq\_jacobi\_mvp\_fft}.
Both functions share the same interface.

%======================================================================
\section{Computational Performance: Old vs.\ New Implementation}
\label{sec:performance}
%======================================================================

This section compares the computational cost of the previous
implementation (Barrowes 2001 block-Toeplitz FFT with \texttt{ray}
multiprocessing) and the current implementation (Goodman 1991 3D FFT
with NumPy/SciPy broadcasting).

%----------------------------------------------------------------------
\subsection{MVP complexity}
\label{sec:mvp-complexity}
%----------------------------------------------------------------------

Both algorithms compute the off-diagonal matrix--vector product via
FFT-based circular convolution, and both have the same asymptotic
complexity per right-hand side (RHS):
\[
  T_{\text{MVP}}^{(1)} = O(N \log N),
\]
where $N = N_x N_y N_z$ is the total cuboid grid size.  For $L$
right-hand sides the total floating-point operation count is
$O(L N \log N)$ in both cases.  The difference lies in execution
efficiency.

%----------------------------------------------------------------------
\subsection{Parallelisation overhead}
\label{sec:parallel-overhead}
%----------------------------------------------------------------------

\begin{table}[htbp]
\centering
\caption{Execution-efficiency comparison (per Krylov iteration =
  two block MVPs for Block-BiCGSTAB).}
\label{tab:perf-compare}
\begin{tabular}{@{}lll@{}}
\toprule
  & Old (\texttt{ray}) & New (broadcasting) \\
\midrule
Parallelism granularity
  & Process-level (1 MVP / worker)
  & Thread-level (FFT + BLAS) \\
Inter-process communication
  & Array serialise/deserialise
  & None (single process) \\
MVP wall time ($L$ RHS)
  & $\sim \dfrac{L}{P}\,T_{\text{MVP}} + T_{\text{ray}}$
  & $\sim \dfrac{L}{W}\,T_{\text{FFT}} \times 2 + T_{\text{matmul}}$ \\
\texttt{ray} overhead $T_{\text{ray}}$
  & $\sim$ 10--100\,ms / MVP
  & 0 \\
Worker startup
  & $\sim$ 1--5\,s (pool init)
  & 0 \\
Cache efficiency
  & Low (each worker holds copy)
  & High (contiguous memory) \\
\bottomrule
\end{tabular}
\end{table}

Here $P$ is the number of \texttt{ray} workers and $W$ is the number
of \texttt{scipy.fft} threads (default: $\texttt{cpu\_count} - 2$).

%----------------------------------------------------------------------
\subsection{Worked example}
\label{sec:perf-example}
%----------------------------------------------------------------------

For $N = 30\,000$, $L = 50$, $P = W = 8$, and 10 Krylov iterations
(20~MVPs):

\begin{table}[htbp]
\centering
\caption{Estimated wall-time breakdown.}
\label{tab:perf-example}
\begin{tabular}{@{}lll@{}}
\toprule
Cost factor & Old (\texttt{ray}) & New (broadcasting) \\
\midrule
Effective MVP calls
  & $20 \times \lceil 50/8 \rceil = 140$
  & 20 (batched) \\
\texttt{ray} IPC overhead
  & $\sim 140 \times 50\,\text{ms} = 7\,\text{s}$
  & 0 \\
Estimated speedup
  & 1$\times$ (baseline)
  & $\sim$ 3--5$\times$ \\
\bottomrule
\end{tabular}
\end{table}

The three main sources of speedup, in decreasing order of impact, are:
\begin{enumerate}
\item \textbf{Removal of \texttt{ray} IPC overhead.}
      Array serialisation/deserialisation is eliminated; the effect
      grows with $N$.
\item \textbf{Batched FFT.}
      All $L$ RHS are processed in a single \texttt{scipy.fft} call;
      the internal thread pool distributes work efficiently.
\item \textbf{Memory locality.}
      Contiguous-array access within a single process improves CPU
      cache hit rates.
\end{enumerate}

%----------------------------------------------------------------------
\subsection{Memory comparison}
\label{sec:memory-compare}
%----------------------------------------------------------------------

\begin{table}[htbp]
\centering
\caption{Peak memory usage.}
\label{tab:memory-compare}
\begin{tabular}{@{}lll@{}}
\toprule
  & Old (\texttt{ray}) & New (Goodman) \\
\midrule
Interaction tensor
  & Copy per worker: $\sim P \times 1152\,N$\,bytes
  & Single copy: $1152\,N$\,bytes \\
Polarisation (during MVP)
  & 1 RHS / worker
  & All $L$ RHS: $768\,L\,N$\,bytes \\
Total peak
  & $\sim P \times 1152\,N$
  & $(1152 + 768\,L) \times N$ \\
\bottomrule
\end{tabular}
\end{table}

For $P = 8$ and moderate $L$, the old implementation uses roughly
$8\times$ more memory for the interaction tensor alone.
The new implementation's memory is dominated by the block polarisation
arrays when $L$ is large.

%======================================================================
\section{Scattering Observables}
\label{sec:observables}
%======================================================================

Once the polarisation $\bm{P}_i$ ($i = 1, \ldots, N$) has been obtained
for each orientation~$\ell$, the following observables are computed.
Throughout this section, $k$ denotes the wavenumber in the medium as
defined in \eqref{eq:wavenumber}.

%----------------------------------------------------------------------
\subsection{Absorption and extinction cross sections}
\label{sec:cross-sections}
%----------------------------------------------------------------------

The absorption cross section is~\cite{Draine1988}
\begin{equation}
  C_{\text{abs}}
    = 4\pi k \sum_{i=1}^{N}
      \operatorname{Im}\!\bigl(\bm{P}_i \cdot \bm{E}_i^*\bigr),
  \label{eq:C-abs}
\end{equation}
where $\bm{E}_i$ is the internal electric field at dipole~$i$,
reconstructed from the polarisation via
\begin{equation}
  \bm{E}_i
    = \frac{4\pi\,\bm{P}_i}{(\epsilon_{r,i} - 1)\,v}.
  \label{eq:E-from-P}
\end{equation}

The extinction cross section is
\begin{equation}
  C_{\text{ext}}
    = 4\pi k \sum_{i=1}^{N}
      \operatorname{Im}\!\bigl(\bm{P}_i \cdot \bm{E}_{\text{inc},i}^*\bigr).
  \label{eq:C-ext}
\end{equation}

The corresponding efficiency factors are
$Q_{\text{abs}} = C_{\text{abs}} / (\pi r_{ve}^2)$ and
$Q_{\text{ext}} = C_{\text{ext}} / (\pi r_{ve}^2)$,
where $r_{ve}$ is the volume-equivalent sphere radius:
\begin{equation}
  r_{ve} = \left(\frac{3\,N\,v}{4\pi}\right)^{\!1/3}.
  \label{eq:r-ve}
\end{equation}

%----------------------------------------------------------------------
\subsection{Forward scattering amplitude (PCAS)}
\label{sec:S-fw}
%----------------------------------------------------------------------

The forward scattering amplitude in the PCAS (Polarimetric Complex
Amplitude of Scattering) framework is computed from the far-field
projection of the dipole polarisations.

For the forward scattering direction $\hat{\bm{u}}_{\text{fw}} = +\hat{\bm{z}}$
in the laboratory frame, define the phase factor for each dipole:
\begin{equation}
  e_{\text{fw},i}
    = \exp\!\bigl(ik\,\hat{\bm{u}}_{\text{fw}} \cdot \bm{r}_i\bigr).
  \label{eq:e-fw}
\end{equation}

The $\theta$- and $\phi$-components of the forward scattering amplitude
correspond to the s- and p-polarisation channels of the CAS-v2
detector, respectively:
$S_{\text{fw},\theta}^{\text{PCAS}} \equiv S_s(0)$ and
$S_{\text{fw},\phi}^{\text{PCAS}} \equiv S_p(0)$.
They are computed as
\begin{align}
  S_{\text{fw},\theta}^{\text{PCAS}}\; [= S_s(0)]
    &= k^2 \sum_{i=1}^{N}
       (\bm{P}_i \cdot \hat{\bm{e}}_{\theta,\text{fw}})\;
       (\sqrt{2}\; e_{\text{fw},i})^*,
  \label{eq:S-fw-theta} \\
  S_{\text{fw},\phi}^{\text{PCAS}}\; [= S_p(0)]
    &= k^2 \sum_{i=1}^{N}
       (\bm{P}_i \cdot \hat{\bm{e}}_{\phi,\text{fw}})\;
       (\sqrt{2}\,i\; e_{\text{fw},i})^*,
  \label{eq:S-fw-phi}
\end{align}
where the factors $\sqrt{2}$ and $\sqrt{2}\,i$ arise from the
circular-polarisation decomposition of the incident field
\eqref{eq:circular-pol}.

%----------------------------------------------------------------------
\subsection{Backward scattering amplitude (OCBS)}
\label{sec:S-bk}
%----------------------------------------------------------------------

For the exact backward scattering direction
$\hat{\bm{u}}_{\text{bk}} = -\hat{\bm{z}}$, define the phase factor
\begin{equation}
  e_{\text{bk},i}
    = \exp\!\bigl(ik\,\hat{\bm{u}}_{\text{bk}} \cdot \bm{r}_i\bigr).
  \label{eq:e-bk}
\end{equation}
Let $\hat{\bm{e}}_{\theta,\text{bk}}$ and
$\hat{\bm{e}}_{\phi,\text{bk}}$ be the polar and azimuthal unit
vectors associated with the backward direction.
In the laboratory frame:
$\hat{\bm{e}}_{\theta,\text{bk}} = -\hat{\bm{x}}$ and
$\hat{\bm{e}}_{\phi,\text{bk}} = +\hat{\bm{y}}$.

The $\theta$- and $\phi$-components of the backward scattering
amplitude are
\begin{align}
  S_{\text{bk},\theta}^{\text{OCBS}}
    &= k^2 \sum_{i=1}^{N}
       (\bm{P}_i \cdot \hat{\bm{e}}_{\theta,\text{bk}})\;
       (\sqrt{2}\; e_{\text{bk},i})^*,
  \label{eq:S-bk-theta} \\
  S_{\text{bk},\phi}^{\text{OCBS}}
    &= k^2 \sum_{i=1}^{N}
       (\bm{P}_i \cdot \hat{\bm{e}}_{\phi,\text{bk}})\;
       (\sqrt{2}\,i\; e_{\text{bk},i})^*,
  \label{eq:S-bk-phi}
\end{align}
where the reference-wave factors $\sqrt{2}$ and $\sqrt{2}\,i$ are the
same as in the forward case.

The OCBS (Opposite-Circular-Basis Scattering) amplitude is then
\begin{equation}
  S_{\text{bk}}^{\text{OCBS}}
    = \frac{-S_{\text{bk},\theta}^{\text{OCBS}}
            + S_{\text{bk},\phi}^{\text{OCBS}}}{\sqrt{2}}.
  \label{eq:S-bk-ocbs}
\end{equation}
More precisely, the $\theta$- and $\phi$-components correspond to
combinations of the Mishchenko amplitude matrix
elements~\cite{Mishchenko2002}:
\begin{equation}
  S_{\text{bk},\theta}^{\text{OCBS}}
    \equiv S_{11}(\pi) + i\,S_{12}(\pi),
  \qquad
  S_{\text{bk},\phi}^{\text{OCBS}}
    \equiv S_{22}(\pi) - i\,S_{21}(\pi).
  \label{eq:S-bk-mishchenko}
\end{equation}
The backscattering theorem $S_{21}(\pi) = -S_{12}(\pi)$ causes the
off-diagonal contributions to cancel when the two components are
combined in \eqref{eq:S-bk-ocbs}, yielding an expression equivalent to
Eq.~(28) of Moteki and Adachi~\cite{Moteki2024}.

%======================================================================
\section{Gaussian Random Ellipsoid (GRE) Shape Model}
\label{sec:gre}
%======================================================================

The particle shape is generated using the Gaussian Random Ellipsoid
(GRE) model of Muinonen \& Pieniluoma~\cite{Muinonen2011}.

%----------------------------------------------------------------------
\subsection{Base ellipsoid}
\label{sec:base-ellipsoid}
%----------------------------------------------------------------------

The base ellipsoid is defined by three semi-axes $a \ge b \ge c$ and
the volume-equivalent radius $r_v$.
Two axis ratios are specified as input: $b/c$ (\texttt{bc\_ratio}) and
$a/b$ (\texttt{ab\_ratio}).
The semi-axis~$c$ is determined by volume conservation:
\begin{equation}
  c = \left(\frac{r_v^3}{(a/b)\,(b/c)^2}\right)^{\!1/3},
  \qquad
  b = c \cdot (b/c),
  \qquad
  a = b \cdot (a/b).
  \label{eq:semi-axes}
\end{equation}

%----------------------------------------------------------------------
\subsection{Surface deformation}
\label{sec:gre-deformation}
%----------------------------------------------------------------------

A Gaussian random field~$s(\theta, \phi)$ is sampled on the surface of
the base ellipsoid with covariance
\begin{equation}
  \operatorname{Cov}(s_i, s_j)
    = \beta^2 \exp\!\left(
        -\frac{\|\bm{r}_i - \bm{r}_j\|^2}{2\,\ell_c^2}
      \right),
  \label{eq:gre-covariance}
\end{equation}
where $\beta$ is the standard deviation of the deformation and
$\ell_c$ is the correlation length, which controls the fineness of
the surface roughness.
In the current version (v0.6.0), $\ell_c = 0.3\,c$ is adopted as an
empirical choice.

The deformed surface position is obtained by displacing each point
on the base ellipsoid along the outward surface normal
$\hat{\bm{n}}_0$:
\begin{equation}
  \bm{r}_{\text{GRE}}(\theta, \phi)
    = \bm{r}_{\text{base}}(\theta, \phi)
      + h_0 \bigl[\exp\!\bigl(s(\theta,\phi)\bigr)
        - \tfrac{1}{2}\beta^2 - 1\bigr]\;
      \hat{\bm{n}}_0(\theta, \phi),
  \label{eq:gre-deformation}
\end{equation}
where $h_0 = c^2 / a$ is a length scale and $\hat{\bm{n}}_0$ is the
unit outward normal to the base ellipsoid.

%----------------------------------------------------------------------
\subsection{Lattice generation and interior detection}
\label{sec:interior-detection}
%----------------------------------------------------------------------

The lattice spacing is determined by the number of dipoles per
wavelength inside the particle (dpl):
\begin{equation}
  d = \frac{\lambda_{\text{particle,min}}}{\mathrm{dpl}}
    = \frac{\lambda_0}{\max_c |m_{p,c}| \times \mathrm{dpl}},
  \label{eq:lattice-spacing}
\end{equation}
where $m_{p,c}$ ($c \in \{x,y,z\}$) is the complex particle
refractive index along axis~$c$ (cf.\ Sec.~\ref{sec:polarisability})
and $\lambda_{\text{particle,min}} = \lambda_0 / \max_c|m_{p,c}|$
is the shortest wavelength inside the particle.
The default value is $\mathrm{dpl} = 17$, which is a typical choice
for DDA simulations of sub-micrometre particles.
A cubic lattice with this spacing is generated within the bounding
box of the deformed surface.
Interior points are identified by a three-axis ray-casting algorithm.

%----------------------------------------------------------------------
\subsection{Volume-preserving lattice rescaling}
\label{sec:volume-rescaling}
%----------------------------------------------------------------------

Because the GRE surface is discretised onto a cubic lattice, the
total volume of the $N$ occupied dipoles, $N d^3$, generally differs
from the target volume $V_{\text{target}} = \tfrac{4}{3}\pi r_{v,\text{base}}^3$
by a few percent.  When parameter sweeps require exact control of the
volume-equivalent radius (e.g.\ an equally-spaced $r_{v,\text{base}}$
grid), this discretisation error is undesirable.

To eliminate it, the lattice spacing is rescaled \emph{after} the
interior detection step.  Given the number of occupied sites~$N$ and
the desired $r_{v,\text{base}}$, the adjusted spacing is
\begin{equation}
  d_{\text{adj}}
    = \left(\frac{V_{\text{target}}}{N}\right)^{\!1/3}
    = \left(\frac{4\pi}{3N}\right)^{\!1/3} r_{v,\text{base}}.
  \label{eq:d-adj}
\end{equation}
All lattice coordinates are then multiplied by the scale factor
$s = d_{\text{adj}} / d$, and $d_{\text{adj}}$ replaces $d$ in every
subsequent computation: the Green's function~\eqref{eq:green-tensor},
the Clausius--Mossotti polarisability~\eqref{eq:cm-polarisability},
and the radiative correction~\eqref{eq:rr-polarisability}.
Because both the inter-dipole distances and the dipole volume
$v = d_{\text{adj}}^3$ change by the same factor, the DDA formulation
remains self-consistent.

This rescaling is performed in the \texttt{Target} constructor
(\texttt{bl\_dda/scatterer.py}), which requires \texttt{r\_v\_base}
as a mandatory argument.

%======================================================================
\section{Spheroid Symmetry and Orientation Averaging}
\label{sec:spheroid}
%======================================================================

For spheroids (axially symmetric particles with $a = b$ and GRE roughness
$\beta_{\mathrm{GRE}} = 0$), the scattering amplitudes possess analytical
constraint relations that enable significant computational savings in
multi-orientation calculations.

%----------------------------------------------------------------------
\subsection{Euler angle mapping}
\label{sec:spheroid-euler}
%----------------------------------------------------------------------

The intrinsic ZYZ Euler angles $(\alpha,\beta,\gamma)$ relate to the
spheroid geometry as follows:
\begin{itemize}
  \item $\beta$: polar angle of the symmetry axis
        (tilt from the laboratory $z$-axis),
  \item $\alpha$: azimuthal angle of the symmetry axis
        (rotation of the tilt plane about $z$),
  \item $\gamma$: rotation about the symmetry axis;
        irrelevant for a spheroid ($a = b$).
\end{itemize}
In the notation of Sec.~\ref{sec:spheroid-constraint},
$\theta \equiv \beta$ and $\phi \equiv \alpha$.

%----------------------------------------------------------------------
\subsection{Constraint relations at forward scattering}
\label{sec:spheroid-constraint}
%----------------------------------------------------------------------

When the symmetry axis lies in the $xz$-plane ($\alpha = 0$), the
cross-polarisation elements of the amplitude scattering matrix vanish:
$S_{12}(\theta,0) = S_{21}(\theta,0) = 0$.
The forward-scattering amplitudes at arbitrary~$\alpha$ are then
determined entirely by the two diagonal elements at $\alpha = 0$:
\begin{align}
  S_s(\theta,\alpha)
    &= \frac{S_{11}(\theta,0) + S_{22}(\theta,0)}{2}
     + \frac{S_{11}(\theta,0) - S_{22}(\theta,0)}{2}\,e^{i2\alpha},
  \label{eq:Ss-constraint} \\
  S_p(\theta,\alpha)
    &= \frac{S_{11}(\theta,0) + S_{22}(\theta,0)}{2}
     - \frac{S_{11}(\theta,0) - S_{22}(\theta,0)}{2}\,e^{i2\alpha},
  \label{eq:Sp-constraint}
\end{align}
where $S_s \equiv S_{11} + iS_{12}$ and $S_p \equiv S_{22} - iS_{21}$
are the PCAS observables (Sec.~\ref{sec:S-fw}).

These imply three verifiable symmetries:
\begin{enumerate}
  \item \textbf{Trace invariance:}\;
    $S_s(\alpha) + S_p(\alpha)
       = S_{11}(\theta,0) + S_{22}(\theta,0) = \text{const}$.
  \item \textbf{Modulus preservation:}\;
    $|S_s(\alpha) - S_p(\alpha)|
       = |S_{11}(\theta,0) - S_{22}(\theta,0)| = \text{const}$.
  \item \textbf{Cross-symmetry:}\;
    $S_s(\theta,\,\alpha+\pi/2) = S_p(\theta,\,\alpha)$.
\end{enumerate}
These relations are verified numerically at machine precision
($\sim 10^{-16}$) in \texttt{test\_spheroid\_symmetry.ipynb}.

%----------------------------------------------------------------------
\subsection{Analytical phi-averaging (labor-saving mode)}
\label{sec:spheroid-saving}
%----------------------------------------------------------------------

For orientation-averaged observables, the azimuthal average
over $\alpha \in [0,2\pi]$ eliminates the oscillatory term
$e^{i2\alpha}$, yielding
\begin{equation}
  \langle S_s(\theta) \rangle_\alpha
  = \langle S_p(\theta) \rangle_\alpha
  = \frac{S_{11}(\theta,0) + S_{22}(\theta,0)}{2}.
  \label{eq:S-phi-avg}
\end{equation}
Furthermore, $C_{\mathrm{ext}}$ and $C_{\mathrm{abs}}$ are also
$\alpha$-invariant for a spheroid, so the values computed at
$\alpha = 0$ are identical to their $\alpha$-averages.

\paragraph{Implementation.}
When \texttt{run\_dda.py} detects a spheroid
(\texttt{ab\_ratio == 1} and \texttt{gre\_beta == 0}),
it activates \emph{spheroid mode}:
\begin{itemize}
  \item The DDA is solved only for $N_\beta$ orientations at
        $\alpha = 0$, $\gamma = 0$, with $\cos\beta$ at the
        $N_\beta$ equal-area midpoints (Sec.~\ref{sec:incident-field}).
  \item The full $N_\alpha \times N_\beta \times N_\gamma$ orientation
        grid is then filled analytically using the constraint relations
        \eqref{eq:Ss-constraint}--\eqref{eq:Sp-constraint}:
        \[
          S_s(\alpha_j, \beta_i)
            = A_i + B_i\,e^{i2\alpha_j}, \qquad
          S_p(\alpha_j, \beta_i)
            = A_i - B_i\,e^{i2\alpha_j},
        \]
        where $A_i = [S_s(\beta_i,0) + S_p(\beta_i,0)]/2$ and
        $B_i = [S_s(\beta_i,0) - S_p(\beta_i,0)]/2$.
  \item For backward scattering, the OCBS amplitude follows the same
        rotational structure:
        $S_{\mathrm{bk}}^{\mathrm{OCBS}}(\alpha_j, \beta_i)
         = S_{\mathrm{bk}}^{\mathrm{OCBS}}(0, \beta_i)\,e^{i2\alpha_j}$.
  \item $C_{\mathrm{ext}}$ and $C_{\mathrm{abs}}$ are $\alpha$-invariant,
        so the $\alpha = 0$ values are broadcast to all grid points.
\end{itemize}
This reduces the number of DDA solves from
$N_\alpha \times N_\beta \times N_\gamma$ to $N_\beta$, while
the output grid retains the full
$N_\alpha \times N_\beta \times N_\gamma$ resolution with
\emph{exact} (not interpolated) values at every grid point.

\paragraph{Applicability.}
These relations hold for any particle whose cross-section in the plane
perpendicular to the symmetry axis is circular ($a = b$), regardless
of the axial ratio $b/c$.  They break down for particles lacking a
continuous rotational symmetry axis (e.g., tri-axial ellipsoids with
$a \neq b$, or particles with GRE roughness $\beta_{\mathrm{GRE}} > 0$).

%======================================================================
\section{Mie Theory Reference}
\label{sec:mie}
%======================================================================

For validation, the DDA solution for a volume-equivalent sphere is
compared to the exact Mie theory~\cite{Bohren1983}.

%----------------------------------------------------------------------
\subsection{Partial-wave expansion}
\label{sec:mie-expansion}
%----------------------------------------------------------------------

For a homogeneous sphere of radius $r_p$, the size parameter is
$x = k\,r_p = 2\pi\,m_m\,r_p / \lambda_0$ and the relative refractive
index is $m_r = m_p / m_m$.
The number of partial-wave terms is
\begin{equation}
  n_{\max} = \lfloor |x + 4\,x^{1/3} + 2| \rfloor.
  \label{eq:nstop}
\end{equation}

The Mie coefficients $a_n$ and $b_n$ are computed from
Riccati--Bessel functions $\psi_n$, $\chi_n$, $\xi_n = \psi_n + i\chi_n$
and the logarithmic derivative $D_n(m_r x)$ via downward
recurrence~\cite{Bohren1983}.

%----------------------------------------------------------------------
\subsection{Efficiencies and scattering amplitudes}
\label{sec:mie-efficiencies}
%----------------------------------------------------------------------

\begin{align}
  Q_{\text{ext}} &= \frac{2}{x^2}\sum_{n=1}^{n_{\max}}
                     (2n+1)\,\operatorname{Re}(a_n + b_n),
  \label{eq:Q-ext-mie} \\
  Q_{\text{sca}} &= \frac{2}{x^2}\sum_{n=1}^{n_{\max}}
                     (2n+1)\,(|a_n|^2 + |b_n|^2),
  \label{eq:Q-sca-mie} \\
  Q_{\text{abs}} &= Q_{\text{ext}} - Q_{\text{sca}}.
  \label{eq:Q-abs-mie}
\end{align}

The angular scattering amplitudes are~\cite{Bohren1983}
\begin{align}
  S_1(\theta)
    &= \sum_{n=1}^{n_{\max}} \frac{2n+1}{n(n+1)}
       \bigl(a_n\,\pi_n(\cos\theta) + b_n\,\tau_n(\cos\theta)\bigr),
  \label{eq:S1-mie} \\
  S_2(\theta)
    &= \sum_{n=1}^{n_{\max}} \frac{2n+1}{n(n+1)}
       \bigl(a_n\,\tau_n(\cos\theta) + b_n\,\pi_n(\cos\theta)\bigr),
  \label{eq:S2-mie}
\end{align}
where $\pi_n$ and $\tau_n$ are the angular functions computed via
upward recurrence.
The forward and backward scattering amplitudes used for validation
are obtained from $S_1$ and $S_2$ at $\theta = 0$ and $\theta = \pi$.

\begin{thebibliography}{99}
\bibitem{Purcell1973}
  E.~M.~Purcell and C.~R.~Pennypacker,
  ``Scattering and absorption of light by nonspherical dielectric grains,''
  \textit{Astrophys.\ J.} \textbf{186}, 705--714 (1973).

\bibitem{Draine1988}
  B.~T.~Draine,
  ``The discrete-dipole approximation and its application to interstellar
  graphite grains,''
  \textit{Astrophys.\ J.} \textbf{333}, 848--872 (1988).

\bibitem{Goodman1991}
  J.~J.~Goodman, B.~T.~Draine, and P.~J.~Flatau,
  ``Application of fast-Fourier-transform techniques to the discrete-dipole
  approximation,''
  \textit{Opt.\ Lett.} \textbf{16}(15), 1198--1200 (1991).

\bibitem{Chaumet2009}
  P.~C.~Chaumet and A.~Rahmani,
  ``Coupled-dipole method for magnetic and negative-refraction materials,''
  \textit{J.\ Quant.\ Spectrosc.\ Radiat.\ Transfer} \textbf{110},
  22--29 (2009).

\bibitem{Tadano2009}
  H.~Tadano, T.~Sakurai, and Y.~Kuramashi,
  ``Block BiCGSTAB($\ell$) for linear systems with multiple right-hand
  sides,''
  \textit{JSIAM Lett.} \textbf{1}, 44--47 (2009).

\bibitem{Gu2016}
  X.-M.~Gu, T.-Z.~Huang, L.~Li, H.-B.~Li, T.~Sogabe, and M.~Clemens,
  ``Block variants of COCG and COCR methods for solving complex symmetric
  linear systems with multiple right-hand sides,''
  arXiv:1611.09813 (2016).

\bibitem{Muinonen2011}
  K.~Muinonen and T.~Pieniluoma,
  ``Light scattering by Gaussian random ellipsoid particles: First results
  with discrete-dipole approximation,''
  \textit{J.\ Quant.\ Spectrosc.\ Radiat.\ Transfer} \textbf{112},
  1747--1752 (2011).

\bibitem{Bohren1983}
  C.~F.~Bohren and D.~R.~Huffman,
  \textit{Absorption and Scattering of Light by Small Particles}
  (Wiley, New York, 1983).

\bibitem{Mishchenko2002}
  M.~I.~Mishchenko, L.~D.~Travis, and A.~A.~Lacis,
  \textit{Scattering, Absorption, and Emission of Light by Small Particles}
  (Cambridge University Press, 2002).

\bibitem{Moteki2024}
  N.~Moteki and K.~Adachi,
  ``Measuring the polarized complex forward-scattering amplitudes of
  single particles in unbounded fluid flow: CAS-v2 protocol,''
  \textit{Opt.\ Express} \textbf{32}(21), 36500--36522 (2024).
  \url{https://doi.org/10.1364/OE.533776}.
\end{thebibliography}

\section*{Acknowledgment}

This document was prepared with the assistance of Claude (Anthropic).
The author assumes full responsibility for the content.

\end{document}
